\begin{abstract}
Speculative decoding improves large language model (LLM) inference throughput by drafting multiple candidate tokens with a smaller model and verifying them with a larger target model. However, reported gains vary substantially across hardware, model pairing, and decoding regime, and consumer-GPU evidence remains limited. We present a controlled empirical study on a single RTX~4090 Laptop using a Qwen2.5-3B-Instruct target and same-family 0.5B/1.5B drafts over GSM8K, a 5-subject MMLU subset, and CNN/DailyMail, with frozen manifests and matched prompts across all runs. We evaluate 12 speculative configurations ($k\in\{4,8,16\}$, deterministic/stochastic regimes, two draft sizes) plus regime-matched baselines, and report wall-clock latency, throughput, acceptance dynamics ($\alpha$, $B_{\text{eff}}$), and task-level quality deltas. On the current RTX4090 CUDA-graph run family, the highest mean speedup is observed for 0.5B, $k{=}4$, deterministic ($S{=}3.0674$). Across the grid, mean speedup decreases with larger draft length ($\bar{S}_{k=4}=2.9164$, $\bar{S}_{k=8}=2.4029$, $\bar{S}_{k=16}=1.8019$), while accepted block size increases and token-level acceptance decreases. Deterministic decoding remains faster in 5 of 6 matched draft--$k$ pairs (mean $\Delta S_{\text{det-stoch}}=0.0568$), with one small exception at 0.5B, $k{=}16$. Quality drift remains modest on MMLU and CNN/DailyMail but is more variable on GSM8K under stochastic settings. We conclude with practical operating recommendations for consumer-GPU deployment and a targeted extension plan to strengthen statistical confidence under coherent paired reruns. Adaptive Cascaded Speculative Decoding (ACSD) is presented as prospective extension work and is not part of the quantitative results reported here.
\end{abstract}